\documentclass[../../../include/open-logic-section]{subfiles}

\begin{document}

\olfileid{sfr}{arith}{rat}
\olsection{从 $\Int$ 到 $\Rat$}

我们刚看到如何运用朴素集合论从自然数构造整数。现在，我们以非常相似的方式从整数构造有理数。
我们初步认识到：
% Part: sets-functions-relations
% Chapter: arithmetization
% Section: rationals

\begin{enumerate}
	\item 每个有理数都可以写成 $\nicefrac{i}{j}$ 的形式，其中 $i$ 和 $j$ 均为整数且 $j$ 非零。
	\item 表达式 $\nicefrac{i}{j}$ 所包含的信息，同样可以用有序对 $\tuple{ i, j}$ 来表示。
\end{enumerate}

最直观的做法是把有理数直接\emph{看作}取自 $\Int \times (\Int \setminus \{0_{\Int}\})$ 的有序对。然而，和之前一样，这样做过于简单，因为我们想要 $\nicefrac{3}{2} = \nicefrac{6}{4}$，却得到 $\tuple{ 3, 2}\neq \tuple{ 6,
4}$。更一般地，我们需要满足以下条件：
\[
	\nicefrac{a}{b} = \nicefrac{c}{d} \text{ 当且仅当 } a \times d = b \times c
\]
为此，我们在 $\Int \times
(\Int \setminus \{0_{\Int}\})$ 上定义一个等价关系如下：
\[
	\tuple{ a, b }\Ratequiv \tuple{ c, d} \text{ 当且仅当 }a \times d = b \times c
\]
我们必须验证这是一个等价关系。这与 $\Intequiv$ 的情形非常类似，我们将其留作练习。

\begin{prob}
证明 $\Ratequiv$ 是一个等价关系。
\end{prob}

借此我们可以说：

\begin{defn}
有理数是（第二个元素非零的）整数对在 $\Ratequiv$ 下的等价类。即，$\Rat =
\equivclass{(\Int \times (\Int\setminus \{0_{\Int}\}))}{\Ratequiv}$。
\end{defn}

与整数的情形类似，我们还需要定义一些基本运算。令 $\equivrep{i,j}{\Ratequiv}$ 表示包含 $\tuple{i, j}$ 作为元素的 $\Ratequiv$ 等价类，我们规定：
\begin{align*}
	\equivrep{a, b}{\Ratequiv} + \equivrep{c, d}{\Ratequiv} &= \equivrep{ad + bc,  bd}{\Ratequiv}\\
	\equivrep{a, b}{\Ratequiv} \times \equivrep{c, d}{\Ratequiv} &= \equivrep{a  c, b d}{\Ratequiv}.
\intertext{为了在这些有理数上定义 $r \leq s$，我们利用 $r \le s$ 当且仅当 $s - r$ 非负这一事实，即 $s - r$ 可以写成 $\nicefrac{i}{j}$ 的形式，其中 $i$ 非负且 $j$ 为正：}
	\equivrep{a, b}{\Ratequiv} \leq \equivrep{c, d}{\Ratequiv} &\text{ 当且仅当 }
	\equivrep{c, d}{\Ratequiv} - \equivrep{a, b}{\Ratequiv} =
	\equivrep{i_{\Int}, j_{\Int}}{\Ratequiv}
\end{align*}
其中 $i \in \Nat$ 且 $0 \neq j \in \Nat$。

接着需要检查这些定义是否如它们所\emph{应当}的那样运作；我们将其留待 \olref[check]{sec}。而它们确实如此！最后，我们需要一种方法将整数\emph{视为}有理数；因此，对于每个 $i \in \Int$，我们规定 $i_{\Rat} =
\equivrep{i, 1_{\Int}}{\Ratequiv}$。我们同样在 \olref[check]{sec} 中验证所有这些是否运作正确。

\begin{prob}
证明对于任意 $i, j \in \Int$，有 $(i + j)_{\Rat} = i_{\Rat}+ j_{\Rat}$ 且 $(i \times j)_{\Rat} =
i_{\Rat} \times j_{\Rat}$ 且 $i \leq j \liff i_{\Rat} \leq j_{\Rat}$。
\end{prob}

\end{document}